\chapter{User Guide}
\label{region-plugin}

\section{Navigation in the region map}
\label{sec-region-navigation}

For each function, the \Region plugin generates a single \emph{region map}:
an overapproximation of a
map dispatching memory locations into regions using alias and shape analysis.
A region map is defined as a graph (acyclic when no pointer is
involved) whose paths start with program variables represented as labels.
\begin{table}[ht!]
\begin{center}
\begin{tabular}{|m{2cm}|m{10cm}|}
\hline
    \centering Region label & Meaning
\\
\hline
{\centering{\includegraphics[width=2cm]{"images/example-region-label.png"}}}
&
A region label (not a region) representing the program variable \texttt{s}.
\\
\hline
\end{tabular}
\end{center}
\vspace{-15px}
\caption{Example of a region label}
\label{table-region-navigation-label}
\vspace{-10px}
\end{table}

\subsection{Region definition}

In any region map, each region is characterized by (at least) its size,
its content, and the memory accesses.
Memory access kinds are: read \texttt{R}, write \texttt{W}, shift
\texttt{S}, and initialization \texttt{I}. Each access is paired with a C type
that defines the size of all memory locations included in the region.

\subsubsection{Primitive regions}

Shown in Table \ref{table-region-navigation-primitive}, primitive regions
are defined as region which only include accesses paired with primitive
C types, which are represented as
\texttt{i} for \texttt{int}, \texttt{s} for \texttt{short},
\texttt{d} for \texttt{double}, \texttt{f} for \texttt{float},
 \texttt{b} for \texttt{bool}, and lastly \texttt{x} for multiple
 heterogeneous accesses.

\begin{table}[ht!]
\begin{center}
\begin{tabular}{|m{4cm}|m{10cm}|}
\hline
    \centering Region & Meaning
\\
\hline
    {\centering{\includegraphics[height=1.1cm]{"images/example-primitive-read.png"}}}
&
    Primitive region with read access(es) of \texttt{int}

    (gray for only primitive read with a unique type).
\\
\hline
    {\centering{\includegraphics[height=1.1cm]{"images/example-primitive-rw.png"}}}
&
    Primitive region with read \& write accesses of \texttt{int}.

    (green for primitive read \& write with a unique type)
\\
\hline
    {\centering{\includegraphics[height=1.1cm]{"images/example-primitive-write.png"}}}
&
    Primitive region with write access(es) of \texttt{int}.

    (pink for only primitive write with a unique type)
\\
\hline
    {\centering{\includegraphics[height=1.4cm]{"images/example-garbled-with-more-details.png"}}}
&
    Primitive region with heterogeneous accesses:

    a 32-bits \texttt{int} and a 32-bits \texttt{float}.

    (red for heterogeneous accesses)
\\
\hline
\end{tabular}
\end{center}
\vspace{-15px}
\caption{Example of primitive regions}
\label{table-region-navigation-primitive}
\vspace{-15px}
\end{table}

\pagebreak
\subsubsection{Pointer regions}

Shown in Table \ref{table-region-navigation-pointer},
pointer regions are regions whose accesses are of C type pointer.
Dereferencing one of the pointer region's pointers (accessing the pointed value)
is interpreted as a pointer read access.
The representation of pointer referencing is shown in the region arc examples
in Table \ref{table-region-navigation-arcs}.

\begin{table}[ht!]
\begin{center}
\begin{tabular}{|m{1.8cm}|m{10cm}|}
\hline
    \centering Region & Meaning
\\
\hline
    {\centering{\includegraphics[height=1.1cm]{"images/example-pointer-read.png"}}}
&
    Pointer region with read access.

    (yellow for only pointer read)
\\
\hline
    {\centering{\includegraphics[height=1.1cm]{"images/example-pointer-rw.png"}}}
&
    Pointer region with read \& write accesses.

    (orange for pointer read \& write)
\\
\hline
    {\centering{\includegraphics[height=1.1cm]{"images/example-pointer-write.png"}}}
&
    Pointer region with write access.

    NB: The addresses contained in this region are never read,
    therefore, no arrow is pointing to another region from it.

    (pink for only pointer write, same as primitive write)
\\
\hline
\end{tabular}
\end{center}
\vspace{-15px}
\caption{Example of pointer regions}
\label{table-region-navigation-pointer}
\vspace{-15px}
\end{table}

\subsubsection{Compound regions}

Shown in Table \ref{table-region-navigation-compound},
compound regions are accessed with a \texttt{struct} C type.
The representation and the shape of the region depends on the accessed
fields of the structure, see Section \ref{sec-region-shape-analysis}.



\begin{table}[ht!]
\begin{center}
\begin{tabular}{|m{2cm}|m{10cm}|}
\hline
    \centering Region & Meaning
\\
\hline
    {\centering{\includegraphics[width=2cm]{"images/example-compound-region.png"}}}
&
    Compound region with at least three fields named \texttt{x}, \texttt{z} and
    \texttt{arr}, where the last one is an array of size \texttt{4}.

    NB: Between the fields \texttt{x} and \texttt{z}, there may be any combination
    of paddings and unaccessed fields of total size 32 bits.

    (no color on structures)
\\
\hline
\end{tabular}
\end{center}
\vspace{-15px}
\caption{Example of a compound region}
\label{table-region-navigation-compound}
\vspace{-15px}
\end{table}


\subsection{Region arcs in the region map}

A region map is defined by its regions and the relations between them.
Those relations are represented as arcs in the region map, as shown in
Table \ref{table-region-navigation-arcs}.

\begin{table}[ht!]
\begin{center}
\begin{tabular}{|m{5cm}|m{10cm}|}
\hline
\centering Region arc & Meaning
\\
\hline
{\centering{\includegraphics[height=1.1cm]{"images/example-labeling.png"}}}
&
The colored primitive region is labelled by the variable named \texttt{\_retres}.
\\\hline
{\centering{\includegraphics[height=1.1cm]{"images/example-labeling-array.png"}}}
&
The colored primitive region is labelled by the variable named \texttt{y}, an array of size 2.
\\
\hline
{\centering{\includegraphics[height=1cm]{"images/example-compound-field.png"}}}
&
The field \texttt{z} on the left is represented by the right colored region.
\\
\hline
{\centering{\includegraphics[height=1.1cm]{"images/example-pointer.png"}}}
&
The right region is pointed by the left pointer region.

It means that all the values in the left pointer region are addresses
included in the right region addresses.
\\
\hline
\end{tabular}
\end{center}
\vspace{-15px}
\caption{Examples of region arcs}
\label{table-region-navigation-arcs}
\vspace{-15px}
\end{table}

\section{Alias and shape analysis}
\label{sec-region-alias-shape-analysis}

The region map analysis is based on the creation of simple region and then
merges each simple region with relevant ones to build more complex regions.


\begin{itemize}
    \item modular
    \item over-approximation
    \item flow-insensitive analysis
    \item field/shape analysis
    \item type access analysis (for heterogenous access analysis)
\end{itemize}

\subsection{Modularity}

The \Region plugin generates a region map for each declared function in the
analysis scope.

\subsection{Flow-insensitive}

For the region map generation, statements in conditional branches are analysed
in the same way as if they were sequential.

This behavior may result into an over-approximation of alias analysis and
unecessary region merges would be avoided when taking into account a deeper
control flow analysis. Yet, the development choice is to provide a simple and
quick region analysis.

\subsection{Shape analysis}
\label{subsec-region-shape-analysis}

The region map is dependent on the machine it is analysed for. For instance,
the same C code will not produce the same region map for two different machdep
where an \texttt{int} is either 32-bits or 16-bits.

\subsubsection{Size analysis}

A region is caracterized by its size which is defined as the biggest size
from which each memory locations' size is a multiple of.

Examples:
\begin{itemize}
    \item Compound region: its size is equal to the overall size of the
    full structure (with possible paddings).
    \item Primitive region with only one C type access: its size is equal to
    the size of the C type.
    \item Pointer region (non-garbled): its size is equal to the size of adresses.
    \item Garbled primitive region: its size is equal to the GCD of all accesses' C type size.
\end{itemize}

\subsubsection{Different sizes merge}

A region is a collection of memory locations. The size of a region is defined
as the greater size of elementary memory blocks necessary to represent all the
memory locations contained in the region. If a memory location has a strictly
greater size than the region size, then this memory location is represented by
an array to correspond to the memory location size.
Each region is caracterized by a size which is definied as the GCD among
all the sizes of the memory locations it includes. A merge of two regions
with different sizes will result into a region that includes a set of
elementary memory blocks of smaller size so that each previous memory location
are now represented as a multiple of the merged elementary memory blocks.

For instance, a region of size 160 bits (e.g. a structure with five 32-bits
\texttt{int}) that is also accessed using a 16-bits \texttt{short} pointer will
degenerate into a garbled primitive region of size 16-bits.

\subsubsection{Compound merge}

\begin{lstlisting}
struct S {
    int a, b, c[4];
};
void job (int k) {
    struct S s;
    s.a = 1;
    s.c[k] = 1;
}\end{lstlisting}

\begin{table}[h!]
\begin{center}
\begin{tabular}{|m{5cm}|m{10cm}|}
\hline
\centering Compatible compound merge & Meaning
\\
\hline
{\centering{\includegraphics[width=5cm]{"images/example-compound-merge-field-a.png"}}}
&
Regions created by analysing the statement \texttt{s.a = 1}.
\\\hline
{\centering{\includegraphics[width=5cm]{"images/example-compound-merge-field-c.png"}}}
&
Regions created by analysing the statement \texttt{s.c[k] = 1}.
\\
\hline \hline
{\centering{\includegraphics[width=5cm]{"images/example-compound-merge-field-a-c.png"}}}
&
Resulting regions from the merge after analysing both statements.

Since field positions and sizes are compatible among themselves,
the resulting region does not collapses into a garbled region.
\\
\hline
\end{tabular}
\end{center}
\vspace{-15px}
\caption{Example of compound merge}
\label{table-compound-merge}
\end{table}
